% Source: https://artofproblemsolving.com/community/c3h2015582_numbers_in_blackboard
% Stable identity: aops-topic-2015582-q01
\textbf{(a)} Replacing $a,b$ by $a+b$ does not change the sum of all numbers on the board. So the final number equals the initial sum:
\[
1+2+\cdots+2017=\dfrac{2017\cdot2018}{2}=\boxed{2035153}.
\]

\textbf{(b)} Note that $ab+a+b+1=(a+1)(b+1)$. Consider the product $P$ of all the numbers $x+1$, for $x$ on the board. Replacing $a,b$ by $c=ab+a+b$ replaces the factors $(a+1)(b+1)$ by $c+1=(a+1)(b+1)$, so $P$ never changes. Initially
\[
P=2\cdot3\cdots2018=2018!,
\]
and at the end $P=z+1$ for the final number $z$. Hence
\[
z=\boxed{2018!-1},
\]
whatever order Ana chooses.
